Respondenternas svar grupperas för varje förbättringsförslag.
Ett bildspel på 17 sidor med scenariobeskrivningar och förslag på meddelandeinnehåll presenterades för respondenterna.
Kontakta författaren för att tillgång detta material.
\paragraph{1.1 Utökat innehåll av JP}
\hspace{-12pt}
\begin{enumerate}
    \item ”Vem” är bra, notis är bra för att slippa sitta och leta ny information.
    \item Bra, mycket av det vi pratade om.
    Något överflödigt med avsändare, däremot bra med telefonnummer.
    Speciellt i utlandet, där kan det också vara intressant att koppla ihop det med bandel/strækning.
    \item Respondenten tyckte att det såg bra ut men tryckte på att det var viktigt att säkerställa frekvens.
    Onödigt med att meddela att man fått ny JP\@.
    \item Respondenten tyckte att det såg bra ut och att det innehöll mycket av de vi pratade om.
    \item Ser absolut rimlig ut, men ser fjärren verkligen det exakta behov av exempelvis ett godståg?
\end{enumerate}
\paragraph{1.2 Blankettförberedelse}
\hspace{-12pt}
\paragraph{1.3 Kontaktuppgifter till tågklarerare}
\hspace{-12pt}
\paragraph{1.4 Trafikinformation}
\hspace{-12pt}
\paragraph{1.5 Tåglägesbild}
\hspace{-12pt}
\paragraph{2.1 Kvittens av JP}
\hspace{-12pt}
\paragraph{2.2 Rapportering balisinformationsfel}
\hspace{-12pt}
\paragraph{2.3 Plattformsbehov, höger/vänster}
\hspace{-12pt}
\paragraph{2.4 Plötsligt stopp, felsökning pågår}
\hspace{-12pt}
\paragraph{2.5 Oplanerat stopp vid stoppsignal}
\hspace{-12pt}
\paragraph{2.6 Plattformsbehov, extra uppehållstid}
\hspace{-12pt}
\paragraph{2.7 GPS-position till SOS Alarm}
\hspace{-12pt}
\paragraph{2.8 Uppdatering, tåglängd}
\hspace{-12pt}
\paragraph{2.9 Samtal önskas}
\hspace{-12pt}